% ===== 2. Theoretical background =====
\section{2. Theoretical background}\label{theoretical-background}

The subject models are decoder-only transformers with causal attention.
The final-layer state at the last position is projected through the
language-model (LM) head onto the vocabulary \(\mathcal{V}\), and
generation feeds each emitted token back as input (Figure
\ref{fig:tokenflow}a, Appendix B.6). The probes read the \emph{prefill}
pass (Figure \ref{fig:tokenflow}b): the single forward pass over the
prompt alone, before any token is generated. Qwen3-1.7B instantiates
this with \(L = 28\) blocks, hidden size \(d = 2048\), and
\(|\mathcal{V}| \approx 152{,}000\) tokens with tied input/output
embeddings (checkpoint and settings in §4).

% ----- 2.1 Confidence metrics: log-probability and entropy -----
\subsection{2.1 Confidence metrics: log-probability and
entropy}\label{confidence-metrics-log-probability-and-entropy}

Two scalar confidence summaries of the generated response recur in §5,
both means over the generated positions. \(\bar\ell\) is the mean
log-probability of the emitted (argmax) tokens, so \(\exp(\bar\ell)\) is
the geometric mean of the per-token probabilities (\(\bar\ell = -0.05\)
corresponds to \(\approx 0.95\) per token). \(\bar H\) is the mean
Shannon entropy of the next-token distribution. All generation is greedy
(\texttt{do\_sample=False}, deterministic), so the emitted token is the
mode. Formal definitions (Eq.~\eqref{eq:lbar}), the worst-step
diagnostic \(\min_t \ell_t\), and the empirical anti-correlation of the
two are in Appendix B.6.

% ----- 2.2 The prefill hidden state as the canonical going-in representation -----
\subsection{2.2 The prefill hidden state as the canonical going-in
representation}\label{the-prefill-hidden-state-as-the-canonical-going-in-representation}

The \emph{prefill} pass over the prompt \(x_{1:T}\) produces the full
hidden-state tensor. I pick off the \textbf{hidden state at every layer
at the last prompt position}, \(\{h^{(\ell)}_T\}_{\ell=0}^{L}\), of
shape \((L+1, d) = (29, 2048)\) per question. Nothing has been generated
yet, so this is the model's representation \emph{immediately before it
commits to the answer} (Figure \ref{fig:tokenflow}b). It summarizes the
whole prompt, sits one linear projection from the first answer token's
logits (Section 4), and avoids the contamination of probing generated
tokens' states (expanded rationale in Appendix B.6).

% ----- 2.3 Why a linear probe (and the specific pipeline I use) -----
\subsection{2.3 Why a linear probe (and the specific pipeline I
use)}\label{why-a-linear-probe-and-the-specific-pipeline-i-use}

Following \citet{alain2017probes}, I fit at each layer \(\ell\) a binary
linear classifier on \(\{h^{(\ell)}_T\}\). Accuracy substantially above
the appropriate baseline implies a hyperplane separates the property at
depth \(\ell\). Whether the model \emph{uses} that direction is a
separate question. The pipeline is
\texttt{StandardScaler\ -\textgreater{}\ PCA(n\_components=16)\ -\textgreater{}\ LogisticRegression(max\_iter=2000,\ C=1.0)}
under 5-fold stratified CV with \texttt{random\_state=0}, reported as
mean and per-fold std of accuracy. Scaler and PCA are fit only on the
training fold, so no test-fold leakage occurs (overfit rationale and
schematic in Appendix B.6/B.7, Figure \ref{fig:pipeline}). Appendix D.1
sweeps \(n_{\mathrm{components}}\) over \(\{4, 8, 16, 32, 48, 64\}\) to
show the headline numbers are not specific to \(16\).

% ----- 2.4 The two confounds and what counts as evidence -----
\subsection{2.4 The two confounds and what counts as
evidence}\label{the-two-confounds-and-what-counts-as-evidence}

\textbf{Confound 1 (cutoff/correctness).} A benchmark of easy
pre-cutoff plus post-cutoff questions makes ``correct'' and
``pre-cutoff'' nearly the same label, so a correctness classifier is in
effect a cutoff classifier (``this question mentions 2025''). ConfabQA
disconfounds with \emph{obscure pre-cutoff} questions, facts in the
training data the model still gets wrong, so pre-cutoff no longer
implies correct (keeping only pre-cutoff questions would also remove the
refusals; Section 5.1). The disconfound test fits the correctness probe
restricted to pre-cutoff items, where accuracy above the pre-cutoff
majority cannot be explained by the cutoff variable.

\textbf{Confound 2 (hidden state/prompt features).} A probe that beats
its majority baseline might still be reading question-text features any
classifier could extract: year, domain, length, keywords. The honest
baseline is a logistic regression on prompt-side features alone (roster
in Section 4, specification in Appendix B.7, schematic in Figure
\ref{fig:baselines}, Appendix B.6). A hidden-state probe is evidence of
model-internal information only to the extent it beats this baseline
outside per-fold noise (\(h_{\rm adds}\) = probe peak \(-\) strongest
baseline, in pp, reported throughout \S\S5--8). This is the control the
probing literature most often omits: it distinguishes ``the hidden state
encodes self-knowledge'' from ``the hidden state encodes the prompt,
which encodes the answer.''
